% Imporditud: tools/import_eksperimendid.py
% Allikas: Eesti füüsikaolümpiaadi eksperimendi ülesannete
%   kogu (Rael Kalda, 2023), ülesanne Ü63, lahendus L63.
\setAuthor{Erkki Tempel}
\setRound{piirkonnavoor}
\setYear{2014}
\setNumber{P E1}
\setDifficulty{2}
\setTopic{staatika}
\setVahendid{tikutops, ühtlase massijaotusega joonlaud, mille mass on teada}
\setPunktid{12}
\setAllikas{Eksperimendi kogu Ü63, lahendus lk 61}
\setLahendusseis{toores}

\prob{Tikutops}
Leidke maksimaalne rõhk, mida laua peal seisev tikke täis tikutops võib lauale avaldada. Tikutopsi katki rebida ei tohi.
\textit{Märkus:} Maksimaalsete punktide saamiseks peab leitud rõhk olema suurem kui 200 Pa.

\hint

\solu
Tikutops avaldab maksimaalset rõhku lauale siis, kui topsi pindala, mis lauale toe-

tub on kõige väiksem. Kõige väiksema pindalaga saab tikutopsi toetada lauale nii,

et lükkad sisemise osa natuke välja ning paned tikutopsi välise kesta servade peale

seisma (3 p).

Toetuspindala leidmine: Tikutopsi väline kest tuleb vajutada kokku (mööda mur-

dejooni), siis saame paksuseks 1 mm ning pikkuseks 4,7 cm, seega pindala on S =

0,1 cm $\cdot$ 4,7 cm = 0,47 cm$^2$ (2 p).

Tikutopsi massi saame leida, kui kasutame joonlauda kangina (2 p). Toetame joon-

laua laua servale ning tikutopsi servapidi teise joonalau otsa ja leiame juhu, kus

joonlaud ulatub maksimaalselt üle laua serva, kuid maha veel ei kuku. Kangi reeg-

list mg $l_1$ + M gl1 = mg $l_2$ saame leida tikutopsi massi, kus $l_1$ on joonlaua lauapealse 2 2 osa pikkus,$l_2$ üle laua oleva osa pikkus, m joonalua mass ning M tikutopsi mass (2

p). Täistikutopsi korral peaks vastus olema vahemikus 6,2-6,6 grammi.

Lauale avaldatud rõhu saame valemist p = F . Vastus tuleb kuskil 1300-1500 Pa (1 S p).

Kordusmõõtmised (2 p).

Märkus: Kui tikutopsi ei panda servade peale seisma, siis saab maksimaalselt 10 p.
\probend
